%!TEX root = main.tex
\setcounter{chapter}{5}
\setcounter{section}{2}
\setcounter{theorem}{0}
\setcounter{equation}{0}

\lectureheader{161}{Calculus I}{Estimation with finite sums}{\textit{Thomas' Calculus} 5.1, \thesection}

\begin{definition}
Given real numbers $a_1,\dots, a_n$, we denote their sum by
\begin{equation*}
\sum_{k=1}^na_k = a_1+\dots + a_n.
\end{equation*}
\end{definition}

\begin{example}
Compute the following.
\begin{enumerate}
\item $\DS\sum_{k=1}^5 k$
\vfill
\item $\DS\sum_{k=1}^3 (-1)^kk$
\vfill
\item $\DS\sum_{k=1}^2 \frac{k}{k+1}$
\vfill
\item $\DS\sum_{k=4}^5 \frac{k^2}{k-1}$
\vfill
\end{enumerate}
\end{example}

\newpage

\begin{theorem}
If $n$ is a positive integer, then
\begin{align}
\sum_{k=1}^nk &= \frac{n(n+1)}{2},\label{sum of consecutive integers}\\
\sum_{k=1}^nk^2 &= \frac{n(n+1)(2n+1)}{6},\\
\sum_{k=1}^nk^3 &= \left(\frac{n(n+1)}{2}\right)^2.
\end{align}
\end{theorem}

\begin{theorem}
Given real numbers $a_1,\dots, a_n$, $b_1,\dots, b_n$, and $c$, we have
\begin{enumerate}
\item $\DS\sum_{k=1}^n\left(a_k+b_k\right) = \sum_{k=1}^na_k + \sum_{k=1}^nb_k$,
\item $\DS\sum_{k=1}^n ca_k = c\sum_{k=1}^na_k$.
\end{enumerate}
\end{theorem}

\begin{example}
Evaluate the following.
\begin{enumerate}[label=(\alph*)]
\item $\DS\sum_{k=1}^n\left(3k-k^2\right)$
\vfill
\item $\DS\sum_{k=1}^n\frac{1}{n+1}$
\vfill
\end{enumerate}
\end{example}

\newpage

\begin{remark}\,
\begin{itemize}
\item Hooke's law says that the force $F$ required to hold a compressed or stretched spring $x$ units from its natural (unstressed) length is proportional to $x$, i.e.,
\begin{equation*}
F=kx
\end{equation*}
for some constant $k$ (measured in force per unit of length).
\item The \textbf{work} done by a \underline{constant} force $F$ on an object as it is displaced $d$ units in a \underline{straight line} in the \underline{direction of the force} is $W=Fd$.
\end{itemize}
\end{remark}

\begin{example}
Consider work required to stretch a spring $1$ meter beyond its natural length if its spring constant is $k=30$ \si{N/m}.
\begin{enumerate}[label=(\alph*)]
\item Why is the answer $W = Fd = \big((30 \si{N/m})\cdot (1 \si{m})\big)\cdot (1 \si{m}) = 30$ \si{N\cdot m} wrong? Is it too big or too little?
\vfill
\item Estimate the work by dividing the 1-meter interval into 2 equal subintervals.
\vspace{5.5in}
\newpage
\item Estimate the work by dividing the 1-meter interval into 4 equal subintervals.
\vfill
\item Estimate the work by dividing the 1-meter interval into 10 equal subintervals.
\vfill
\end{enumerate}
\end{example}

\newpage

\begin{definition}
A \textbf{partition} of the interval $[a,b]$ is a finite set of numbers $P=\{x_0,x_1,\dots, x_n\}$ so that
\begin{equation*}
a=x_0<x_1<\dots < x_n=b.
\end{equation*}
\end{definition}
\begin{definition}
Let $f$ be a positive, continuous function on $[a,b]$, let $P=\{x_0,\dots, x_n\}$ be a partition of $[a,b]$, and let $\Delta_j=x_j-x_{j-1}$ denote the length of the $j$th subinterval.
\begin{itemize}
\item The \textbf{upper sum} approximation to the area under the graph of $f$ is
\begin{equation*}
\sum_{j=1}^n M_j\Delta_j,
\end{equation*}
where $M_j=\max\big\{f(x) : x\in [x_{j-1}, x_j]\big\}$.
\item The \textbf{lower sum} approximation to the area under the graph of $f$ is
\begin{equation*}
\sum_{j=1}^n m_j\Delta_j,
\end{equation*}
where $m_j=\min\big\{f(x) : x\in [x_{j-1}, x_j]\big\}$.
\item The \textbf{midpoint rule} approximation to the area under the graph of $f$ is
\begin{equation*}
\sum_{j=1}^nf(c_j)\Delta_j,
\end{equation*}
where $c_j=(x_j+x_{j-1})/2$ is the midpoint of the interval $[x_{j-1}, x_j]$.
\end{itemize}
\end{definition}

\begin{example}
Estimate the area under under the curve $y=1-x^2$ over the interval $[0,2]$ as follows.
\begin{enumerate}[label=(\alph*)]
\item Use a lower sum with 4 rectangles of equal width.
\vfill
\newpage
\item Use an upper sum with 4 rectangles of equal width.
\vfill
\item Use the midpoint rule with 4 rectangles of equal width.
\vfill
\end{enumerate}
\end{example}
